% !TeX program = xelatex
\documentclass[a4paper,11pt]{article}
\usepackage[margin=25mm]{geometry}
\usepackage{fontspec}
\usepackage{kotex}
\setmainfont{Latin Modern Roman}
\IfFontExistsTF{Malgun Gothic}{%
  \setmainhangulfont{Malgun Gothic}%
}{%
  \IfFontExistsTF{Noto Serif CJK KR}{%
    \setmainhangulfont{Noto Serif CJK KR}%
  }{%
    \IfFontExistsTF{UnBatang}{\setmainhangulfont{UnBatang}}{}%
  }%
}
\usepackage{amsmath,amssymb,amsthm}
\usepackage[unicode,hidelinks]{hyperref}
\hypersetup{
  pdftitle={격자점 색칠과 모든 방향의 단색 직사각형},
  pdfauthor={Math Project Lean},
  pdfsubject={최소 높이의 완전한 증명과 Lean 형식 검증}
}
\linespread{1.16}
\setlength{\parskip}{0.3em}
\setlength{\parindent}{1em}
\setlength{\emergencystretch}{2em}
\allowdisplaybreaks[1]
\theoremstyle{definition}
\newtheorem{theorem}{정리}
\newtheorem{lemma}{보조정리}
\renewcommand{\proofname}{증명}
\newcommand{\TT}{T}
\newcommand{\OO}{\mathrm{O}}
\newcommand{\SSp}{\mathrm{S}}

\title{격자점 색칠과 모든 방향의 단색 직사각형\\[0.3em]
  \large 최소 높이의 구성적 증명}
\author{Math Project Lean}
\date{}

\begin{document}
\maketitle
\vspace{-1.2em}

\noindent\textbf{문제.}
양의 정수 $n$에 대하여, 가로 $n+1$열과 세로 $\ell$행의 격자점을
$n$가지 색으로 칠한다. 어떻게 칠하더라도 같은 색인 네 점이 직사각형의
꼭짓점이 되게 하는 $\ell$의 최솟값을 구하여라.
직사각형의 변은 좌표축과 평행하지 않아도 된다.

이 문서에서 격자는 가로와 세로의 간격이 같은 \textbf{정사각 격자}이다.
따라서 점들의 좌표를 $(x,y)$, $0\le x\le n$, $0\le y<\ell$인 정수로
잡는다. 직사각형은 넓이가 양수인 것을 뜻하며, 주어진 색을 모두 사용해야
한다는 조건은 없다.

\begin{theorem}\label{thm:main}
문제에서 요구하는 최소 높이는
\[
  \boxed{\ell_{\min}=n\binom{n+1}{2}+1
  =\frac{n^2(n+1)}{2}+1}
\]
이다.
\end{theorem}

증명은 두 부분으로 나뉜다. 먼저 비둘기집 원리로 이 높이에서 축평행
단색 직사각형이 반드시 생김을 보인다. 이어 한 행 적은 높이에 대한
색칠을 구성하고, 기울어진 직사각형까지 모두 배제한다.

\section{비둘기집 원리에 의한 상한}

한 행에는 $n+1$개의 점이 있으므로 같은 색인 두 점이 있다.
각 행에서 그러한 쌍 하나를 골라 \textbf{두 열의 위치와 그 색}을 기록한다.
서로 다른 기록은 모두 $n\binom{n+1}{2}$가지이다.

따라서 행이 $n\binom{n+1}{2}+1$개이면 두 행의 기록이 같다.
그 두 행과 두 열이 만나는 네 점은 같은 색이고 축평행 직사각형을 이룬다.
그러므로 정리의 값은 최소 높이의 상한이다.

\section{한 행 적은 높이의 색칠 구성}

다음과 같이 삼각수와 블록의 크기를 정의한다.
\[
  T_b=\binom b2=\frac{b(b-1)}2,
  \qquad M=\binom{n+1}{2}=T_{n+1}.
\]
색 번호는 $0,1,\ldots,n-1$로 표시한다. $nM$개의 행을 $M$행씩 $n$개
블록으로 나누고, 블록 번호를 $s=0,1,\ldots,n-1$로 둔다.
각 블록의 행은 다음과 같이 정확히 한 번씩 표현된다.
\[
  y=sM+T_b+a,\qquad 1\le b\le n,\quad 0\le a<b.
\]
이는 $[T_b,T_b+b-1]=[T_b,T_{b+1}-1]$이므로 이 구간들이
$0,1,\ldots,M-1$을 순서대로 분할하기 때문이다.

이 행의 열 $x$에 다음 색을 칠한다.
\begin{equation}\label{eq:coloring}
 C(x,sM+T_b+a)\equiv s+
 \begin{cases}
   x,&0\le x<b,\\
   a,&x=b,\\
   x-1,&b<x\le n
 \end{cases}
 \pmod n.
\end{equation}
즉 $0,1,\ldots,n-1$을 나열하다가 열 $b$에 색 $a$를 하나 더 끼워
넣은 뒤, 모든 색에 블록 번호 $s$를 더한다.
아래의 행 배열은 $y$가 증가하는 순서로 적었다.

예를 들어 $n=3$의 첫 블록은
\[
\begin{array}{c|cccc}
 y&x=0&x=1&x=2&x=3\\ \hline
 0&0&0&1&2\\
 1&0&1&0&2\\
 2&0&1&1&2\\
 3&0&1&2&0\\
 4&0&1&2&1\\
 5&0&1&2&2
\end{array}
\]
이다. 여기에 각 색에 $1$, $2$를 더한 블록을 차례로 붙인다.
덧셈을 $\bmod 3$으로 계산하면 총 $18$행의 색칠을 얻는다.

\subsection*{축평행 단색 직사각형의 배제}

행 $(s,a,b)$에서 같은 색인 유일한 열 쌍은 $(a,b)$이고, 그 색은
$s+a\pmod n$이다. 고정된 열 쌍 $(a,b)$에 대하여 $s$를 바꾸면
각 색이 정확히 한 번씩 나온다. 따라서 어떤 \textbf{열 쌍과 색}도
서로 다른 두 행에서 반복되지 않는다.
이 색칠에는 축평행 단색 직사각형이 없다.

남은 목표는 기울어진 단색 직사각형도 없음을 보이는 것이다.
제\ref{sec:height}절부터 제\ref{sec:wrap}절까지는 $n\ge3$을 가정한다.
$n=1,2$는 제\ref{sec:small}절에서 별도로 확인한다.

\section{기울어진 직사각형의 높이 제한}\label{sec:height}

기울어진 격자 직사각형의 최하단 꼭짓점을 $B$라 하자.
$B$에서 출발하는 두 변의 벡터는
\[
  (p,q),\qquad (-r,t),\qquad p,q,r,t\in\mathbb Z_{>0}
\]
로 쓸 수 있다. 두 변이 직교하고 격자의 가로 폭이 $n$이므로
\begin{equation}\label{eq:rectangle-parameters}
  qt=pr,\qquad p+r\le n.
\end{equation}
직사각형의 높이는 $q+t$이다. $(q-1)(t-1)\ge0$을 이용하면
\begin{equation}\label{eq:height}
  q+t\le qt+1=pr+1
  \le\left\lfloor\frac{n^2}{4}\right\rfloor+1=:H.
\end{equation}

이제 $n\ge3$이면 다음 부등식들이 성립한다.
\begin{equation}\label{eq:height-comparisons}
  H\le T_n<M-1<M,
  \qquad H<M-n+1=T_n+1.
\end{equation}
첫 번째 부등식은 $n=3$일 때 $H=T_n=3$으로 확인된다.
$n\ge4$일 때에는
\[
  H\le\frac{n^2}{4}+1\le\frac{n(n-1)}2=T_n
\]
을 쓰면 된다. 나머지는 $M=T_n+n$에서 바로 따른다.

특히 높이가 $M$보다 작으므로 기울어진 직사각형의 꼭짓점들은
하나의 블록 또는 인접한 두 블록에만 속한다.
이 두 경우를 나누어 살펴보자.

\section{한 블록 안에서의 색의 모양}

블록 $s$에서 하나의 전역 색을 고정하고, 이 색을 $s+k\pmod n$으로
나타내는 $0\le k<n$을 잡는다. 블록 내부의 행 좌표를 $u=y-sM$으로
쓰면, 이 색의 점들은 다음 두 종류로 정확히 나뉜다.

\begin{itemize}
\item \textbf{보통점 $\OO$.} 각 행에 하나씩 있으며 그 열은
\begin{equation}\label{eq:ordinary-one}
  x=\begin{cases}
    k+1,&0\le u<T_{k+1},\\
    k,&T_{k+1}\le u<M.
  \end{cases}
\end{equation}
\item \textbf{특별점 $\SSp$.} 다음 점들이 하나씩 더 있다.
\begin{equation}\label{eq:special-one}
  (x,u)=(b,T_b+k),\qquad b=k+1,k+2,\ldots,n.
\end{equation}
\end{itemize}

이를 확인하려면 식~\eqref{eq:coloring}에서 색 $k$가 나오는 열을
추적하면 된다. 삽입 위치 $b\le k$인 행에서는 기존 색 $k$가
열 $k+1$로 밀린다. $b>k$인 행에서는 열 $k$에 그대로 있다.
추가로 삽입된 색이 $k$인 경우는 $a=k<b$이며, 바로
식~\eqref{eq:special-one}의 점들이다.

\subsection*{한 블록 안의 기울어진 직사각형 배제}

보통점은 위로 갈수록 열이 감소하거나 일정하다. 서로 다른 두 보통점의
연결선은 음의 기울기를 갖거나 수직이며, 수평일 수 없다.
따라서 세 보통점은 직각삼각형을 이루지 못한다.

서로 다른 두 특별점은 열이 다르고, 연결선의 기울기는
\begin{equation}\label{eq:within-slope}
  \frac{T_b-T_a}{b-a}=\frac{a+b-1}{2}>0
  \qquad(a<b)
\end{equation}
이다. 세 특별점 역시 직각삼각형을 이루지 못한다.
직사각형의 꼭짓점 중 아무 세 개를 택해도 직각삼각형이 되므로,
단색 직사각형이 존재한다면 보통점 두 개와 특별점 두 개로 이루어져야 한다.

모든 보통점의 열은 $k+1$ 이하이고 모든 특별점의 열은 $k+1$ 이상이다.
특별점 두 개의 열은 서로 다르므로
\[
  x(\OO_1)+x(\OO_2)<x(\SSp_1)+x(\SSp_2).
\]
직사각형의 두 대각선은 중점이 같으므로 이 두 쌍은 대각선일 수 없다.
따라서 각각 마주보는 변이어야 한다. 그러나 보통점 쌍의 연결선은
음의 기울기를 갖거나 수직이고, 특별점 쌍의 연결선은 양의 기울기를
가지므로 서로 평행하지 않다. 이것도 불가능하다.

\section{인접 블록의 일반 경계를 넘는 경우}\label{sec:ordinary-boundary}

블록 번호가 $s$에서 $s+1$로 바뀌면, 고정된 전역 색의 기본 색 번호는
$k$에서 $k-1\pmod n$으로 바뀐다.
먼저 $1\le k\le n-1$인 경우를 다룬다.
아래 블록의 시작 높이를 $0$으로 평행이동하면 두 블록은 각각
$0\le y<M$, $M\le y<2M$이다.

\subsection{보통점의 폭과 좌우 순서}

두 블록을 합친 보통점의 열은
\begin{equation}\label{eq:ordinary-two}
 x=\begin{cases}
  k+1,&y<T_{k+1},\\
  k,&T_{k+1}\le y<M+T_k,\\
  k-1,&M+T_k\le y.
 \end{cases}
\end{equation}
두 번의 전환 사이의 간격은
\[
  (M+T_k)-T_{k+1}=M-k\ge M-n+1>H.
\]
따라서 높이 $H$ 이하인 구간에서 만나는 보통점은 기껏해야 인접한
두 열에만 놓인다. 또한 열은 위로 갈수록 감소하거나 일정하므로,
앞 절과 마찬가지로 세 보통점이 직각삼각형을 이룰 수 없다.

특별점들은 아래 블록과 위 블록에서 각각
\begin{equation}\label{eq:special-two}
  (b,T_b+k)\quad(k+1\le b\le n),\qquad
  (a,M+T_a+k-1)\quad(k\le a\le n)
\end{equation}
이다. 높이 $H$ 이하인 같은 구간에서는 \textbf{모든 보통점이 모든
특별점보다 왼쪽 또는 같은 열에 놓인다.}
식~\eqref{eq:ordinary-two}와 \eqref{eq:special-two}에서 이 순서가
어긋날 유일한 가능성은 아래 블록의 열 $k+1$ 보통점과 위 블록의
열 $k$ 특별점이다. 하지만 이들의 높이 차는 최소
\[
 (M+T_k+k-1)-(T_{k+1}-1)=M>H
\]
이므로 같은 구간에 나타날 수 없다.

\subsection{특별점 연결선의 기울기}

$P=M-1$로 놓자. 두 블록의 같은 열에 놓인 특별점들의 높이 차는
$P>H$이다. 따라서 하나의 기울어진 직사각형에 포함될 특별점들은
서로 다른 열에 놓인다.

같은 블록의 두 특별점 사이의 기울기는 식~\eqref{eq:within-slope}에
의하여 $1$ 이상이며, $1$인 경우는 두 열이 $1,2$일 때뿐이다.
다른 블록의 특별점을 연결할 때, 아래 점의 열을 $b$, 위 점의 열을
$a$라 쓰자. $a>b$이면 그 기울기는
\[
  \frac{P+T_a-T_b}{a-b}
  =\frac{P}{a-b}+\frac{a+b-1}{2}>1.
\]
$a<b$이면 기울기는 음수이고 그 절댓값은
\begin{equation}\label{eq:negative-slope}
  \frac{P-(T_b-T_a)}{b-a}\ge1.
\end{equation}
실제로 $1\le a<b\le n$이므로
\[
  (T_b-T_a)+(b-a)=T_{b+1}-T_{a+1}
  \le T_{n+1}-T_2=M-1=P.
\]
삼각수의 엄격한 증가에 의하여 등호는 $a=1$, $b=n$일 때만 성립한다.

\subsection{특별점이 세 개 이상인 경우의 배제}

특별점 사이의 연결선은 수직도 수평도 아니고, 모든 기울기의 절댓값은
$1$ 이상이다. 두 연결선이 직교하려면 기울기의 곱이 $-1$이어야
하므로, 두 기울기는 반드시 $+1$과 $-1$이다.

기울기가 $+1$인 특별점 쌍은 같은 블록의 열 $1,2$에 놓인다.
아래 블록의 특별점 열은 $k+1\ge2$이므로 이 쌍은 위 블록에 있고,
$k=1$이어야 한다. 기울기가 $-1$인 쌍은 앞의 등호 조건에 따라
위 블록의 열 $1$과 아래 블록의 열 $n$에 놓인다.
따라서 직각을 이루는 세 특별점의 유일한 형태는 어떤 $Y$에 대하여
\begin{equation}\label{eq:exceptional}
  (1,Y),\qquad (2,Y+1),\qquad(n,Y-n+1)
\end{equation}
이다. 직각은 $(1,Y)$에서 이루어진다.
이를 직사각형으로 완성하는 네 번째 꼭짓점은
\[
  (2,Y+1)+(n,Y-n+1)-(1,Y)=(n+1,Y-n+2)
\]
이므로 격자의 열 범위 $0\le x\le n$을 벗어난다.
따라서 특별점 세 개를 포함하는 단색 직사각형은 없다.

\subsection{보통점 두 개와 특별점 두 개인 경우의 배제}

세 보통점도, 세 특별점도 직사각형에 포함될 수 없으므로 남은 경우는
보통점 두 개와 특별점 두 개이다. 앞서 보인 좌우 순서와 특별점들의
서로 다른 열에 의하여
\[
  x(\OO_1)+x(\OO_2)<x(\SSp_1)+x(\SSp_2).
\]
따라서 두 쌍은 대각선일 수 없고, 마주보는 변이어야 한다.
마주보는 변은 평행할 뿐 아니라 길이도 같으므로 두 변의 벡터는
부호를 제외하면 같다.

보통점 두 개가 같은 열에 있다면 특별점 두 개도 같은 열에 있어야
하므로 불가능하다. 남은 경우 보통점의 열 차는 $1$이고 기울기는
음수이다. 따라서 특별점 쌍 역시 열 차가 $1$이고 기울기가 음수이다.
이는 서로 다른 블록의 특별점이어야 한다.

아래 특별점의 열이 $b+1$, 위 특별점의 열이 $b$라면
$1\le b\le n-1$이고, 두 점의 세로 차는
\[
  D=M-1+T_b-T_{b+1}=M-1-b\ge M-n.
\]
기울어진 직사각형에서는 다른 방향의 변도 세로 성분의 절댓값이
양의 정수이다. 직사각형 전체 높이는 두 방향의 세로 성분 절댓값의
합이므로 적어도
\[
  D+1\ge M-n+1>H
\]
이다. 이는 식~\eqref{eq:height}의 높이 상한과 모순이다.
이로써 일반 경계를 넘는 경우도 모두 배제되었다.

\section{색 번호가 순환하는 경계를 넘는 경우}\label{sec:wrap}

마지막으로 기본 색 번호가 아래 블록에서 $0$, 위 블록에서 $n-1$인
경우를 보자. 아래 블록의 시작 높이를 다시 $0$으로 잡는다.
직사각형이 경계를 넘으므로 최하단 높이는 $M-1$ 이하이다.
따라서 위 블록에 속한 꼭짓점의 내부 행 좌표 $u=y-M$에 대해서는
다음이 성립한다.
\[
  0\le u\le H-1<T_n
\]
기본 색 $n-1$의 점은 식~\eqref{eq:ordinary-one}과
\eqref{eq:special-one}에 의하여 이 범위에서 모두 열 $n$에 있다.

특히 직사각형의 최상단 꼭짓점이 열 $n$에 놓인다.
하지만 기울어진 직사각형에서는 최우측 꼭짓점이 최상단 꼭짓점보다
엄격히 오른쪽에 있다. 제\ref{sec:height}절의 벡터 표기로 쓰면
최우측 꼭짓점은 $B+(p,q)$, 최상단 꼭짓점은 $B+(p-r,q+t)$이고,
두 점의 가로 차는 $r>0$이다.
따라서 최우측 꼭짓점의 열이 $n$을 초과하게 되어 모순이다.

\section{작은 경우와 최소성의 결론}\label{sec:small}

$n=1$일 때에는 $2$열 $1$행에 직사각형이 없고, $2$열 $2$행에서는
네 점이 모두 같은 색이다. 따라서 답은 $2$이다.

$n=2$일 때 식~\eqref{eq:coloring}의 구성은 다음 여섯 행이다.
\[
\begin{array}{c|ccc}
 y&x=0&x=1&x=2\\ \hline
 0&0&0&1\\
 1&0&1&0\\
 2&0&1&1\\
 3&1&1&0\\
 4&1&0&1\\
 5&1&0&0
\end{array}
\]
축평행 단색 직사각형이 없다는 증명은 이미 주었다.
기울어진 경우에는 식~\eqref{eq:rectangle-parameters}에서
$p+r\le2$, $qt=pr$이므로 $p=q=r=t=1$이다.
따라서 가능한 형태는
\[
  (0,y),\quad(1,y-1),\quad(2,y),\quad(1,y+1)
\]
의 마름모뿐이다. 양 끝 열이 같은 색인 행은 $y=1,4$뿐이다.
$y=1$일 때 이웃한 두 행의 가운데 색은 각각 $0,1$이고,
$y=4$일 때에는 각각 $1,0$이다. 어느 쪽도 네 점이 모두 같은 색은
아니므로 이 구성에는 기울어진 단색 직사각형도 없다.

이제 모든 $n\ge1$에 대하여 $nM$행의 반례 색칠을 얻었다.
이 색칠의 처음 $\ell$행만 남기면 $\ell\le nM$인 모든 높이에서도 반례가 된다.
반면 $nM+1$행에서는 제1절에 의하여 단색 직사각형이 반드시 생긴다.
따라서
\[
  \boxed{\ell_{\min}=nM+1
  =n\binom{n+1}{2}+1=\frac{n^2(n+1)}2+1}
\]
이며, 정리~\ref{thm:main}의 증명이 끝났다. \hfill$\square$

\clearpage
\section{Lean 형식 증명과의 대응}

이 문제의 Lean 4 형식 증명은 이 문서와 같은 폴더를 기준으로
\path{lean/RectangleColoring/}에 있다.
\path{Basic.lean}은 정수 격자, 색칠, 비퇴화 직사각형과 최소성의
의미를 정의한다. 직사각형은 대각선 중점의 일치와 인접한 두 변의
직교로 정의하므로 모든 방향이 포함된다.

최종 정리는 \path{lean/RectangleColoring/Main.lean}의
\texttt{RectangleColoring.minimum\_height}이다.
\begin{verbatim}
theorem minimum_height (n : Nat) (hn : 0 < n) :
    IsMinimum n (n ^ 2 * (n + 1) / 2 + 1)
\end{verbatim}
여기서 \texttt{IsMinimum n l}은 높이 \texttt{l}의 모든 색칠이 단색
직사각형을 포함하고, 더 작은 모든 높이에서는 그렇지 않음을 뜻한다.
프로젝트의 증명에는 \texttt{sorry}, \texttt{admit}, 별도로 추가한
공리, \texttt{native\_decide}가 없다. 최종 정리의 공리 의존성은
Lean/Mathlib의 표준 논리 기반인 \texttt{propext},
\texttt{Classical.choice}, \texttt{Quot.sound}이다.

\path{verify.py}는 명시적인 색칠에 대한 별도의 유한 범위 검산 도구이다.
Lean의 일반 정리는 이 Python 검산에 의존하지 않는다.
설치된 도구 버전과 빌드 명령은 저장소의 \path{README.md}에 정리되어 있다.

\end{document}
